\section{Conclusion and Future Work}
\label{sec:conclusion}

In this work, we introduced \textbf{ELATI} (\textbf{E}early-\textbf{L}ayer \textbf{A}daptive \textbf{T}ask \textbf{I}dentification), a training-free, parameter-free ``one-pass'' dynamic model merging framework designed to resolve the systems latency bottleneck of prior penultimate-layer routing systems.
By shifting the routing decision to Layer~2, ELATI theoretically eliminates the expensive throw-away first pass of the network backbone, achieving near-single-pass serving latency.
To perform routing without semantic heads, ELATI introduces Early-Layer Representative Mapping (ELRM) to project intermediate activations against unsupervised task centroids computed from a hyper-sparse 64-sample calibration split.
On heterogeneous streams, Downstream-Only Micro-Batch Homogenization (DO-MBH) dynamically groups activations on-the-fly and interpolates only downstream model parameters.

Our empirical evaluation on a high-fidelity \textbf{Hierarchical 14-Layer Sandbox} across 10 random seeds demonstrates that ELATI achieves a highly robust Joint Mean accuracy of \textbf{56.89\% $\pm$ 1.66\%} (a \textbf{+8.62\%} absolute gain over static uniform merging) on simulated task manifolds.
Furthermore, simulated end-to-end forward propagation benchmarks demonstrate a genuine \textbf{1.40$\times$ physical CPU speedup}, reducing end-to-end latency of our simulated deep network from \textbf{36.90~ms} (PFSR) to \textbf{26.43~ms} on CPU, validating the theoretical operations reduction.
Additionally, CPU micro-benchmarks of the routing projection step show that ELATI reduces latency from \textbf{1.31~ms} to \textbf{0.39~ms} in fully vectorized mode (a \textbf{3.33$\times$ speedup} on CPU) and from \textbf{171.05~ms} to \textbf{63.48~ms} in sequential loop mode (a \textbf{2.69$\times$ speedup}) compared to the penultimate PFSR baseline, driven by reducing theoretical complexity from $O(B \cdot K \cdot C \cdot D)$ to $O(B \cdot K \cdot D)$.
Crucially, ELATI remains extremely robust to overfitting under data-scarcity, outperforming classical linear classifiers on noisy tasks, and maintains robust positive utility under representational entanglement sweeps.

\paragraph{Future Directions and Real-World Systems Scaling}
To scale ELATI to large-scale serving of massive parameter models (e.g., LLaMA-7B or 70B), future work must address the physical memory bandwidth bottleneck of dynamic weight materialization.
For a model with $P$ parameters and $K$ experts, dynamic interpolation requires reading $P \cdot K$ expert parameters and writing $P$ merged parameters to VRAM, which is highly memory-bandwidth-limited. For a 7B-parameter model in FP16, this requires reading and writing $\approx 14$ GB of weight tensors, which can introduce millisecond-level systems overhead on hardware.
To resolve this physical bottleneck, we propose two hardware-native serving strategies:
\begin{enumerate}
    \item \textbf{Low-Rank Downstream Arithmetic (PEFT Serving):} Instead of materializing full merged weights $W_{\text{merged}} = W_{\text{base}} + \sum \bar{\alpha}_k V_k^{(l)}$ in VRAM, serving engines (such as Punica \cite{chen23punica} or S-LoRA \cite{sheng23slora}) perform the base forward pass and low-rank adapter projections on-the-fly without full-weight materialization:
    \begin{equation}
        y = x W_{\text{base}} + \sum_{k=1}^K \bar{\alpha}_k \left( (x A_k) B_k^T \right)
    \end{equation}
    Because $A_k, B_k$ are low-rank matrices ($r = 8$), they are over $1000\times$ smaller than $W_{\text{base}}$, completely bypassing the full-weight materialization memory bandwidth limit.
    \item \textbf{Subspace Caching and Layer-wise Asynchronous Copying:} Overlapping weight interpolation with token execution by pre-fetching downstream layer weights into GPU registers asynchronously while early layers are being computed in Pass 1.
\end{enumerate}

Importantly, when deploying ELATI on physical GPU systems, practitioners must consider the specific hardware interconnect topography and memory architecture. On discrete GPU nodes (e.g., standard PCIe-attached accelerators), the overhead of copying weights between Host RAM and GPU VRAM via PCIe lanes can introduce non-trivial latency, making asynchronous pre-fetching and double-buffering schemes critical. Conversely, on unified memory architectures (e.g., Apple Silicon M-series or NVIDIA Grace Hopper Superchips), where the CPU and GPU share a high-bandwidth unified physical memory pool with a single physical address space, Host-to-Device copying is completely bypassed. On such unified architectures, dynamic weight interpolation and parameter ensembling are virtualized and incur near zero-overhead memory traffic, significantly enhancing the practical systems speedup of ELATI.

Evaluating these dynamic weight caching, low-rank arithmetic, and dynamic offloading strategies on actual physical GPUs is a critical and exciting next step. Specifically, our upcoming hardware roadmap includes deploying ELATI within a custom fork of S-LoRA or vLLM to profile actual HBM-to-SRAM weight loading latencies, register allocation, and kernel execution overheads on physical GPU architectures (e.g., NVIDIA H100 or L4). This physical implementation will serve as a key milestone to bridge our theoretical and simulated findings with production-grade serving systems.
By offering an optimized trade-off between task performance and systems complexity, ELATI establishes a highly transparent, deployable, and low-cost path for dynamic parameter merging on resource-constrained consumer hardware.
